\section{The \instname\ data reduction flow.}
\label{sec:reduction_flow}


The overall data flow of the \instname\ pipeline is displayed in Figure \ref{fig:reduction_cascade}.

 
\begin{figure}[h!]
\centering
\includegraphics[width=0.9\textwidth]{../instrument_tutorial/figures/reduction_cascade.jpg}
\caption{The data reduction cascade of the \instname\ workflow.}
\label{fig:reduction_cascade}
\end{figure}



The reduction cascade is organized in tasks, which represent
well-defined steps in the process. Tasks can be grouped inside
sub-workflows.  Each task runs a recipe; the detailed description of
the algorithms, input, outputs and recipe parameters used in each
recipe are available in the pipeline manual. Here, we present only the
description of most important features.

The \wkfname\  \edps\ workflow is designed to execute the tasks that
deliver the final reduced data cube for each dataset. It can be either
the product of a single exposure, or the combination of multiple
exposures. Only calibrations needed by the selected the scientific
exposures are processed.

It is possible to set \edps\ to perform the data reduction until a
certain step of the reduction chain (e.g. to reduce only standar
stars, or only flat fields).  This is done by specifying the desired
tasks in the field {\bf Select reduction target} of the {\bf Raw Data}
tab.

The reduction steps of the \wkfname\ workflow are listed below. Before
starting the reduction, the parameters of the recipes associated to
each task can be configured by pressing the button

\includegraphics[width=0.8cm]{../common/figures/configure_dataset.jpg}
close to each dataset configuration.
See for more info on the configuration editor \ref{sec:configuration_editor}




\subsection{Generate Master Bias}
This step is carried in the task {\bf bias}, which runs the recipe  {\tt xsh\_mbias}.

Produces a master bias for UVB/VIS arms. NIR frames do not use bias frames.

\bigskip
{\bf Customization}
\bigskip

Recipe parameters:
\begin{itemize}
\item Choose stacking method (average, median) with {\tt stack-method}.
\item Adjust sigma-clipping with {\tt klow} and {\tt khigh}.
\end{itemize}

\subsection{Generate Master Dark}
This step is carried in the task {\bf dark}, which runs the recipe  {\tt xsh\_mdark}.

Produces a master dark.
Optional for UVB/VIS (dark current negligible).
For NIR it is only needed for stare observations; 
generally not used in nodding/offset where ON--OFF subtraction removes the dark. 

\bigskip
{\bf Customization}
\bigskip

Recipe parameters:
\begin{itemize}
\item Enable the use of darks for UVB/VIS data reduction with {\tt use\_optical\_dark}.
\item Change {\tt stack-method} and {\tt klow} / {\tt khigh} thresholds to adjust stacking behaviour, e.g. for better cosmic-ray rejection.
\end{itemize}

\subsection{Fit Orders}

Computes initial guesses for the wavelength solution and order positions.
Performed in two steps, described hereafter.

\subsubsection{Instrument Model Prediction}

Recipe: {\tt xsh\_predict}

Uses the instrument physical model and information about 
ambient conditions during the observations 
(e.g. atmospheric  pressure, temperature, instrument setting)
to predict line positions and a first dispersion solution.

\bigskip
{\bf Customization}
\bigskip

Workflow parameters:
\begin{itemize}
\item Select between physical-model (recommended) and polynomial mode. 
See {\tt configure\_reduction.md} for more information.
\end{itemize}

\subsubsection{Order Tracing (Determining Order Geometry)}

Recipe: {\tt xsh\_orderpos}

Uses pinhole order-definition frames (ORDERDEF\_*) 
to trace the location and curvature of each echelle order on the detector. 
Produces order tables that later help rectification and wavelength calibration.

\bigskip
{\bf Customization}
\bigskip

Recipe parameters:
\begin{itemize}
\item Tuning of detection thresholds can help when continuum levels are low or orders are partially vignetted.
\end{itemize}

\subsection{Generate Master Flat}

Recipe: {\tt xsh\_mflat}

Produces a master flat,
refined version of the table from {\tt xsh\_orderpos}
and a bad pixel map for each arm. 
The flat defines the detector response and 
updates the true geometry of the orders from slit illumination.
For IFU flat fields the edges of the slices are traced.

\bigskip
{\bf Customization}
\bigskip

Recipe parameters:
\begin{itemize}
\item Adjust bad-pixel handling (via {\tt decode-bp}) if master flats saturate or raise quality control errors.
\item Change {\tt stack-method} and {\tt klow} / {\tt khigh} thresholds to adjust stacking behaviour, e.g. for better cosmic-ray rejection.
\end{itemize}

\subsection{Wavelength Calibration}

The wavelength calibration creates the wavelength and spatial resampling solutions 
and computes the arc-line tilts and instrumental resolution.
It is done in two steps:

\subsubsection{2D Mapping}

Recipe: {\tt xsh\_2dmap}

Determines the two-dimensional wavelength solution needed to resample the orders.
Creates {\tt WAVE\_TAB\_2D\_ARM}, 
which provides a full mapping from detector coordinates to wavelength+slit coordinates, 
and {\tt SPECTRAL\_FORMAT\_TAB\_ARM},
which defines, for each order, the wavelength range, pixel boundaries, predicted order edge traces,
and spatial-to-spectral coordinate conversion.

\bigskip
{\bf Customization}
\bigskip

Workflow parameters:
\begin{itemize}
\item Select between physical-model (recommended) and polynomial mode.
\end{itemize}

Recipe parameters:
\begin{itemize}
\item Adjust line-detection thresholds for weak arcs via {\tt detectarclines-fit-win-hsize} and {\tt detectarclines-search-win-hsize}. 
These parameters have to be small enough not to include a doublet but large
enough to be able to detect and fit the line.
\end{itemize}

\subsubsection{Wavelength Calibration}

Recipe: {\tt xsh\_wavecal}

Computes arc lines tilt and resolving power.

\bigskip
{\bf Customization}
\bigskip

Workflow parameters:
\begin{itemize}
\item Select between physical-model (recommended) and polynomial mode.
\end{itemize}

\subsection{Flexure Compensation}

Recipe: {\tt xsh\_flexcomp}

Refines the wavelength solution to account for flexure, 
especially when arcs are not taken at the same rotator angle as science.

\bigskip
{\bf Customization}
\bigskip

Workflow parameters:
\begin{itemize}
\item Select between physical-model (recommended) and polynomial mode.
\end{itemize}

\subsection{Flat Strategy}

This allows the user to choose between two strategies 
to select a flat field for the flux calibrator.
The default strategy is to use the same flat as for the science observation.
The alternative is to use the flat field selected by the rules, 
i.e., those taken closest in time to the flux calibrator.

\bigskip
{\bf Customization}
\bigskip

Workflow parameters:
\begin{itemize}
\item By default, the {\tt use\_flat} parameter is set to {\tt science}, 
meaning the flats used for science frames are also applied to standard stars.
Set it to {\tt standard} to use the flats taken closest in time 
to the standard-star observations.
\end{itemize}

\subsection{Instrument Response and Efficiency}

Recipes: {\tt xsh\_respon\_slit\_stare}, {\tt xsh\_respon\_slit\_offset} and {\tt xsh\_respon\_slit\_nod}

Use standard stars to derive the per-order and merged instrument response 
(mapping detector counts to physical flux),
the blaze correction and the telescope + instrument + detector efficiency.
Used for flux calibration of science exposures. 
The pipeline does not create response curves for IFU data, and they are therefore
not flux-calibrated.

\subsection{Science Reduction SLIT}

Recipes: {\tt xsh\_scired\_slit\_stare}, {\tt xsh\_scired\_slit\_offset} and {\tt xsh\_scired\_slit\_nod}

Perform the data reduction: 
\begin{itemize}
\item Prepare the science frame (bias/dark/inter-order background correction, flat-fielding).
\item Rectify orders using the model and wavelength solution.
\item Perform sky subtraction (mode-dependent).
\item Localize the object and extract {\tt ORDER1D} and {\tt MERGE1D} products.
\end{itemize}

When several exposures are fed together, 
the recipes stack them (mean/median) before extraction, 
and you get one combined 2D and 1D spectrum per run, not one per exposure.
To obtain one spectrum per exposure (or per AB pair in nodding), 
you must run {\tt xsh\_scired\_slit\_*} separately on each exposure or nod pair.

\bigskip
{\bf Customization}
\bigskip

Workflow parameters:
\begin{itemize}
\item By default, {\tt telluric\_correction\_mode}={\tt standard} derives 
the atmospheric parameters from the telluric standard.
Set it to {\tt science} to derive them directly from the science frame, 
or to {\tt none} to disable telluric correction.
\item Select between physical-model (recommended) and polynomial mode.
\end{itemize}

Recipe parameters:
\begin{itemize}
\item Change {\tt stack-method} and {\tt klow} / {\tt khigh} thresholds to adjust stacking behaviour.
\item Sky modelling: 
Use {\tt sky-method} to select the method. 
In the NIR, {\tt BSPLINE2} provides the best residuals,
{\tt BSPLINE1} is faster and acceptable for UVB/VIS, 
and {\tt MEDIAN} is the fatest but may leave residuals.
\item Sky region selection:
use parameters {\tt sky-position1}, {\tt sky-hheight1}, {\tt sky-position2}, and {\tt sky-hheight2} to manually select sky zones.
Required when object not centered, multiple objects on slit, or strong NIR gradient.
\item Spectroscopic extraction:
Standard extraction 
({\tt localize-method}={\tt MANUAL})
recommended for faint sources.
Automatic detection
({\tt localize-method}={\tt AUTO})
is usually fine for bright sources.
\end{itemize}

